\begin{proposition}[DP26: a finite numerical saving is not excluded]\label{DP26}
For the ramps of DP17, let $A=p^3$ and $w_i=1$. Then
\[
 R_2=p^2-1=(p-1)(p+1),\qquad
 R_3=p^2+p-2=(p-1)(p+2),
\]
so $(p+1)R_3=(p+2)R_2$. At $p=11$ this gives
\[
 12\cdot130=13\cdot120,
 \qquad 24\cdot130=26\cdot120=3120.
\]
The second equality is compatible with even lift counts. It does not construct a theta chain or assert that its other costs are unchanged.
\end{proposition}
\begin{proof}
Only $i=1,2$ can contribute positively: their terms are $p^2-1$ and $p-1$. The $i=3$ term is zero, and all $i\ge4$ give $p^{3-i}<1$. Height two retains only the first positive term; height three retains both. This proves the formulas. The source-shaped finite cutoff is \path{range}(\lfloor\log_p A\rfloor+2)=\{0,1,2,3,4\}$, giving the same values. The term $i=0$ is excluded by both divisibility-index conventions, and every weight satisfies $w_i\ge1$.
\end{proof}
